

% --- resumo em português --- (Obrigatório para TCC2)
\begin{resumo}
	A gestão operacional em canteiros de obras ainda depende de processos manuais de coleta e comunicação de informações, o que dificulta o monitoramento das atividades, aumenta o tempo de resposta e favorece a ocorrência de erros. Nesse contexto, os Gêmeos Digitais destacam-se como uma tecnologia capaz de integrar dados em tempo real e fornecer uma representação virtual intuitiva do processo construtivo, apoiando a tomada de decisão. Entretanto, falhas operacionais decorrentes de erros humanos durante a montagem de sistemas construtivos, como as formas metálicas do sistema construtivo de paredes de concreto, ainda representam um desafio significativo para o controle de qualidade, exigindo mecanismos capazes de detectar precocemente não conformidades e minimizar seus impactos na produtividade e na execução da obra.
	
	Este trabalho propõe uma metodologia baseada em Gêmeo Digital que integra Simulação Computacional Híbrida e diagnose de falhas para o monitoramento de processos de montagem e desmontagem de fôrmas. A modelagem híbrida combina Modelagem Baseada em Agentes (ABM), responsável por representar o comportamento dos trabalhadores, e Simulação a Eventos Discretos (DES), utilizada para modelar o fluxo operacional do sistema. A diagnose de falhas fundamenta-se na teoria de Sistemas a Eventos Discretos (SED), permitindo inferir a ocorrência de falhas não observáveis a partir da sequência de eventos observáveis do processo, por meio da construção de um Autômato Diagnosticador.
	
	A simulação foi desenvolvida no software AnyLogic com base em dados de obras reais, enquanto a integração com a plataforma NVIDIA Omniverse proporciona visualização tridimensional fotorrealista e sincronizada do ambiente virtual. A integração entre a Simulação Computacional Híbrida, a diagnose de falhas baseada em Sistemas a Eventos Discretos e a visualização imersiva compõe um Gêmeo Digital capaz de apoiar a detecção precoce de falhas operacionais, alertar sobre pontos críticos do canteiro de obras e auxiliar a tomada de decisão do gestor.
	
	\vspace{\onelineskip}
	\noindent
	\textbf{Palavras-chave:} Gêmeos Digitais; Simulação Computacional Híbrida; Modelagem Baseada em Agentes; Simulação a Eventos Discretos; Autômatos; Sistemas a Eventos Discretos; NVIDIA Omniverse; Anylogic; Diagnose de Falhas; Construção Civil.
\end{resumo}

% --- resumo em inglês --- (Obrigatório para TCC2)
\begin{resumo}[Abstract]
	\begin{otherlanguage*}{english}
		Construction site operational management still relies heavily on manual processes for data collection and information exchange, making it difficult to monitor activities, increasing response times, and contributing to the occurrence of errors. In this context, Digital Twins have emerged as a promising technology capable of integrating real-time data and providing an intuitive virtual representation of construction processes, thereby supporting managerial decision-making. However, operational failures caused by human errors during the assembly of construction systems, such as the steel formwork used in concrete wall construction, remain a significant challenge for quality control, requiring mechanisms capable of detecting nonconformities at an early stage and minimizing their impact on productivity and project execution.
		
		This work proposes a Digital Twin-based methodology that integrates Hybrid Simulation and fault diagnosis for monitoring the assembly and disassembly processes of formwork systems. The hybrid modeling approach combines Agent-Based Modeling (ABM), which represents workers' behavior, and Discrete Event Simulation (DES), which models the operational workflow of the system. Fault diagnosis is based on the theory of Discrete Event Systems (DES), enabling the inference of unobservable faults from sequences of observable events through the construction of a Diagnoser Automaton.
		
		The simulation was developed using the AnyLogic software based on data collected from real construction projects, while integration with the NVIDIA Omniverse platform provides synchronized photorealistic three-dimensional visualization of the virtual environment. The integration of Hybrid Simulation, fault diagnosis based on Discrete Event Systems, and immersive visualization composes a Digital Twin capable of supporting the early detection of operational failures, identifying critical points within the construction site, and assisting managers in the decision-making process.
		
		\vspace{\onelineskip}
		\noindent
		\textbf{Keywords}: Digital Twins; Hybrid Computational Simulation; Agent-Based Modeling; Discrete Event Simulation; Automata; Discrete Event Systems; NVIDIA Omniverse; AnyLogic; Fault Diagnosis; Construction Industry..
	\end{otherlanguage*}
\end{resumo}


